\documentclass[11pt]{article}

\usepackage[a4paper,margin=2.4cm]{geometry}
\usepackage[utf8]{inputenc}
\usepackage[T1]{fontenc}
\usepackage{lmodern}
\usepackage{graphicx}
\usepackage{booktabs}
\usepackage{tabularx}
\sloppy
\usepackage{xcolor}
\usepackage{hyperref}
\usepackage{enumitem}
\usepackage{amsmath}
\usepackage{tcolorbox}
\usepackage{caption}
\usepackage{titlesec}
\usepackage{fancyhdr}

\graphicspath{{outputs/r4/}{outputs/r4/req3_eta_comparison/}}

\definecolor{inkblue}{HTML}{1B3A5C}
\definecolor{amberflag}{HTML}{9A5B15}
\definecolor{qagreen}{HTML}{2E5E3E}
\definecolor{softgrey}{HTML}{6B7680}

\titleformat{\section}{\large\bfseries\color{inkblue}}{\thesection}{0.6em}{}
\titleformat{\subsection}{\normalsize\bfseries\color{inkblue}}{\thesubsection}{0.5em}{}

\hypersetup{
  colorlinks=true, linkcolor=inkblue, citecolor=inkblue, urlcolor=inkblue,
  pdftitle={Requirement 4 -- Study Notes}, pdfauthor={OLA Project}
}

\pagestyle{fancy}
\fancyhf{}
\renewcommand{\headrulewidth}{0.3pt}
\lhead{\small\color{softgrey} OLA -- Requirement 4}
\rhead{\small\color{softgrey} Note di studio}
\rfoot{\small\color{softgrey} \thepage}

\newtcolorbox{flagbox}{colback=amberflag!8, colframe=amberflag, boxrule=0.6pt, arc=2pt, left=8pt, right=8pt, top=6pt, bottom=6pt}
\newtcolorbox{qabox}[1]{colback=qagreen!6, colframe=qagreen, boxrule=0.6pt, arc=2pt, left=8pt, right=8pt, top=6pt, bottom=6pt, title={\small\bfseries #1}, fonttitle=\color{white}, colbacktitle=qagreen}
\newtcolorbox{keybox}{colback=inkblue!6, colframe=inkblue, boxrule=0.6pt, arc=2pt, left=8pt, right=8pt, top=6pt, bottom=6pt}

\setlist[itemize]{leftmargin=1.4em, itemsep=2pt, topsep=3pt}

\title{\color{inkblue}\bfseries Requirement 4 --- Slightly Non-Stationary Environment\\[4pt]
\large\normalfont\color{softgrey} Note di studio -- Online Learning Applications, Prof.\ Castiglioni}
\author{}
\date{}

\begin{document}
\maketitle
\thispagestyle{fancy}
\vspace{-1.2cm}

\begin{keybox}
\textbf{Come usare queste note.} Non sono lo script della presentazione (per quello vedi l'artifact web) -- sono pensate per essere lette e capite a fondo, incluse le domande che il professore potrebbe fare. Ogni sezione con immagine cita il file esatto nella cartella \texttt{outputs/r4/}.
\end{keybox}

\tableofcontents

%============================================================
\section{Inquadramento generale}
%============================================================

Requirement 4 riusa \emph{esattamente} lo stesso problema di Requirement 2 e 3: $N=4$ campagne, budget condiviso $B=1600$, orizzonte $T=10\,000$, grafo dei conflitti $\{(0,1),(2,3)\}$, valori $v = [0.8, 0.8, 0.9, 0.9]$, griglia di bid \texttt{linspace(0,1,11)} e feedback completo (osserviamo $m_t$, la miglior offerta avversaria, alla fine di ogni round).

L'unica cosa che cambia rispetto a Req3 è \textbf{il tipo di non-stazionarietà}: Req3 usa \texttt{mode='drift'} (cambio continuo, ogni round), Req4 usa \texttt{mode='shocks'} (cambio discreto, a blocchi). È la stessa classe di ambiente (\texttt{AdversarialMultiCampaignEnv}), solo riparametrizzata.

\begin{table}[h]
\centering
\small
\begin{tabularx}{\textwidth}{@{}l X X@{}}
\toprule
& \textbf{Req3 -- drift} & \textbf{Req4 -- shocks} \\
\midrule
Tipo di cambiamento & continuo, ogni round & discreto, a blocchi \\
Meccanismo & $m_t \sim \text{Beta}(\alpha_t,\beta_t)$, media sinusoidale & 5 blocchi da 2000 round, regime i.i.d.\ per blocco \\
Prevedibilità & la direzione del cambiamento è liscia & nessun preavviso al confine di blocco \\
Benchmark naturale & $\text{OPT}^A$ (miglior fissa nel tempo) & clairvoyant piecewise (per-blocco) \\
\bottomrule
\end{tabularx}
\caption{Confronto tra i due tipi di non-stazionarietà usati in Req3 e Req4.}
\end{table}

%============================================================
\section{L'ambiente: shock piecewise-stazionari}
%============================================================

L'orizzonte $T=10\,000$ è diviso in $5$ blocchi da $2000$ round (\texttt{block\_size=2000}, \texttt{n\_regimes=5}). Prima dell'inizio della simulazione vengono estratti $5$ regimi $\text{Beta}(\alpha_i,\beta_i)$ con media $\mu_i \sim \text{Uniform}(0.1, 0.9)$; a ogni blocco viene assegnato uno di questi $5$ regimi \textbf{scelto uniformemente a caso} (quindi un regime può ripetersi, non è garantito un giro completo su tutti e 5).

Dentro un blocco, $m_t$ è i.i.d.\ dal regime assegnato. Al confine tra blocchi, la distribuzione cambia istantaneamente e senza segnale. L'agente non conosce:
\begin{itemize}
\item dove sono i confini dei blocchi;
\item quanti regimi distinti esistono;
\item se un regime già visto si ripeterà.
\end{itemize}

\begin{flagbox}
\textbf{Limite noto, da citare se chiesto:} poiché i regimi sono riassegnati ai blocchi in modo i.i.d., non c'è garanzia che i 5 blocchi coprano tutti i 5 regimi distinti -- in alcune run un regime può comparire due volte e un altro zero volte. Non invalida i risultati (il benchmark piecewise usa comunque le distribuzioni vere di ciascun blocco), ma è un dettaglio implementativo onesto da conoscere se il professore lo chiede.
\end{flagbox}

%============================================================
\section{Le tre strategie}
%============================================================

\subsection{Sliding-Window Combinatorial-UCB}
Estende il \texttt{CombinatorialUCBAgent} di Req2 mantenendo le statistiche solo sugli ultimi $W$ round (finestra temporale, non per-braccio). Con $W=500$ (un quarto della lunghezza di un blocco): abbastanza corto da smaltire i dati stantii nei primi $\sim$500 round dopo uno shock, abbastanza lungo da lasciare $\sim$1500 round di sfruttamento stabile.

\subsection{CUSUM Combinatorial-UCB}
Stessa base, ma con un rilevatore di cambiamento Page (1954) CUSUM per ogni cella (campagna, bid): quando la statistica CUSUM supera una soglia $h$, la cella viene resettata (trattata come mai esplorata). Parametri usati: $h=8.0$, $\varepsilon=0.02$, $U_T=4$ (numero atteso di cambi di regime). Con probabilità $\alpha=\sqrt{U_T\log(T/U_T)/T}$ gioca un'azione casuale come rete di sicurezza contro derive lente che CUSUM potrebbe non cogliere subito.

\subsection{Primal-Dual (Req3, riusato)}
\`E \emph{lo stesso identico agente} di Requirement 3 (Hedge per campagna + discesa duale OGD su $\lambda$), \textbf{nessuna modifica strutturale} -- solo i tassi di apprendimento sono stati ritarati per Req4: \texttt{hedge\_eta=0.16}, \texttt{ogd\_eta=0.003}, con \texttt{budget\_pacing=True} (target duale adattivo $\rho_t = B_{\text{rimanente}}/T_{\text{rimanente}}$).

\begin{qabox}{D: perché Primal-Dual non è stato ridisegnato per Req4?}
Perché il punto pedagogico del requirement è proprio verificare cosa succede quando si \emph{riusa senza modifiche strutturali} un algoritmo con garanzia di best-of-both-worlds in un ambiente per cui quella garanzia non è stata pensata. Ridisegnarlo avrebbe nascosto il fenomeno interessante.
\end{qabox}

%============================================================
\section{Metodologia dei benchmark}
%============================================================

Come in Req3, un clairvoyant dinamico (che conosce ogni $m_t$ realizzato) gonfia il regret di un termine lineare in $T$ per la disuguaglianza di Jensen -- nessun learner può avere regret sub-lineare contro di esso in un'asta first-price. Usiamo quindi tre benchmark, in ordine di importanza:

\begin{table}[h]
\centering
\small
\begin{tabularx}{\textwidth}{@{}l X l@{}}
\toprule
\textbf{Livello} & \textbf{Cosa conosce} & \textbf{Funzione} \\
\midrule
\textbf{Primario} & confini dei blocchi + distribuzione vera per blocco (non i singoli $m_t$) & \small\texttt{compute\_piecewise\_}\newline\small\texttt{expected\_clairvoyant} \\
Secondario ($\text{OPT}^A$) & miglior distribuzione fissa nel tempo (continuità con Req3) & \small\texttt{compute\_clairvoyant\_}\newline\small\texttt{multi} \\
Riferimento (prophet) & ogni $m_t$ realizzato & \small\texttt{compute\_clairvoyant\_}\newline\small\texttt{dynamic\_multi} \\
\bottomrule
\end{tabularx}
\caption{Gerarchia dei tre benchmark usati in Req4.}
\end{table}

Il clairvoyant piecewise è il target naturale per SW-UCB e CUSUM-UCB: ha esattamente le garanzie teoriche di letteratura (es.\ Garivier \& Moulines, 2011) contro cui questi algoritmi sono progettati.

%============================================================
\section{Risultati principali}
%============================================================

\begin{table}[h]
\centering
\small
\begin{tabular}{@{}lrrrr@{}}
\toprule
\textbf{Agente} & \textbf{Piecewise} & $\textbf{OPT}^A$ & \textbf{Prophet} & \textbf{Costo finale / B} \\
\midrule
Primal-Dual (Req4-tuned) & \textbf{1426.1} & 626.8 & 2699.7 & 1456.4 / 1600 (91.0\%) \\
CUSUM Combinatorial-UCB & 1430.0 & 630.7 & 2703.6 & 1599.3 / 1600 (99.96\%) \\
Sliding-Window Comb.-UCB & 1947.6 & 1148.3 & 3221.2 & 1599.2 / 1600 (99.95\%) \\
\bottomrule
\end{tabular}
\caption{Risultati finali medi su 20 trial, $T=10\,000$ (numeri dal notebook eseguito).}
\end{table}

\begin{figure}[h]
\centering
\includegraphics[width=0.85\textwidth]{req4_regret_piecewise.png}
\caption{Regret cumulativo contro il clairvoyant piecewise (benchmark primario). Primal-Dual e CUSUM sono sostanzialmente appaiati; Sliding-Window è nettamente peggiore.}
\end{figure}

\begin{figure}[h]
\centering
\includegraphics[width=0.85\textwidth]{req4_budget.png}
\caption{Costo cumulativo dei tre agenti contro la linea del budget. CUSUM e Sliding-Window esauriscono quasi tutto il budget; Primal-Dual si ferma al 91\%.}
\end{figure}

\textbf{Lettura critica.} Il ranking sul regret non racconta tutta la storia. Guardando la spesa di budget, Primal-Dual è chiaramente diverso dagli altri due: CUSUM e Sliding-Window hanno un meccanismo esplicito di reset innescato da un cambio di regime rilevato; Primal-Dual no. Questo è il punto di partenza della Sezione~\ref{sec:tuning}.

%============================================================
\section{\`E un problema di taratura?}
\label{sec:tuning}
%============================================================

Per verificare se il sotto-utilizzo del budget da parte di Primal-Dual fosse dovuto a una taratura non ottimale (ereditata da Req3, pensata per il drift continuo), abbiamo confrontato \textbf{lo stesso agente, sullo stesso ambiente shocks}, con due configurazioni di learning rate:

\begin{table}[h]
\centering
\footnotesize
\begin{tabularx}{\textwidth}{@{}Xrrrrr@{}}
\toprule
\textbf{Configurazione} & \textbf{Piecewise} & $\textbf{OPT}^A$ & \textbf{Prophet} & \textbf{Costo / B} & $\boldsymbol{\lambda}$ \textbf{fin.} \\
\midrule
Req3 eta (baseline, non ritarato) & 1777.4 & 978.0 & 3051.0 & 1193.9/1600 (74.6\%) & 0.575 \\
Req4 tuned (\texttt{hedge=0.16}, \texttt{ogd=0.003}) & \textbf{1426.1} & 626.8 & 2699.7 & 1456.4/1600 (91.0\%) & 1.018 \\
\bottomrule
\end{tabularx}
\caption{Confronto Primal-Dual: tassi di Req3 (non ritarati per lo shock discreto) vs tassi ritarati per Req4.}
\end{table}

\begin{figure}[h]
\centering
\includegraphics[width=0.85\textwidth]{comparison_budget.png}
\caption{Costo cumulativo: Req3 eta (arancio, baseline) vs Req4 tuned (blu). Nessuna delle due curve tocca mai la linea rossa del budget.}
\end{figure}

\begin{figure}[h]
\centering
\includegraphics[width=0.85\textwidth]{comparison_regret_piecewise.png}
\caption{Regret piecewise nel tempo: il retuning riduce sensibilmente il regret lungo tutto l'orizzonte, non solo al termine.}
\end{figure}

\begin{figure}[h]
\centering
\includegraphics[width=0.75\textwidth]{comparison_lambda.png}
\caption{Traiettoria di $\lambda_t$: in entrambe le configurazioni $\lambda$ impiega centinaia di round a risalire dopo ogni shock -- il retuning accorcia il ritardo, non lo elimina.}
\end{figure}

\textbf{Conclusione empirica.} Il retuning \emph{aiuta}: il regret scende da 1777 a 1426 (arriva a pareggiare CUSUM) e la spesa sale dal 75\% al 91\%. Ma \textbf{nessuna delle due configurazioni si avvicina al 99.9\% di CUSUM e Sliding-Window}. Aumentare il tasso di apprendimento del primale (Hedge) accelera l'adattamento \emph{dentro} ogni blocco; non risolve il ritardo con cui $\lambda$ reagisce \emph{subito dopo} uno shock -- quel ritardo è strutturale al meccanismo di discesa duale OGD, che non riceve mai un segnale esplicito "il regime è appena cambiato".

\begin{qabox}{D: perché non avete cercato tassi ancora più aggressivi?}
Un tasso di apprendimento di Hedge troppo alto causa un'instabilità numerica reale: i pesi moltiplicativi $w \mathrel{*}= \exp(-\eta \cdot \text{loss})$ vanno in underflow a zero esatto, e la normalizzazione $w/\sum w$ diventa $0/0=\text{NaN}$. C'è quindi un limite fisico oltre al quale l'algoritmo smette semplicemente di funzionare, non solo di migliorare.
\end{qabox}

%============================================================
\section{Perché è strutturale: la teoria in breve}
%============================================================

\begin{itemize}
\item \textbf{Hedge (il primale di Primal-Dual)} ha una garanzia di \emph{regret statico/esterno}: converge, in aspettativa, alla miglior distribuzione \emph{fissa nel tempo} -- esattamente la definizione di $\text{OPT}^A$. Non ha alcun meccanismo di \emph{dimenticanza} o di rilevamento di cambio regime incorporato nell'algoritmo.
\item \textbf{OGD sul duale $\lambda$} insegue un target $\rho_t$ con un passo di gradiente proporzionale allo scarto tra spesa realizzata e target -- una risposta \emph{graduale}, adatta a un target che si muove con continuità (Req3, drift), non a un target che salta discretamente (Req4, shocks).
\item \textbf{SW-UCB e CUSUM-UCB}, al contrario, hanno un meccanismo di \emph{tracking regret} esplicito: la finestra scorrevole scarta attivamente i dati vecchi; CUSUM resetta esplicitamente alla detection di uno shift. Sono progettati, in letteratura, proprio per ambienti piecewise-stazionari.
\end{itemize}

\begin{keybox}
\textbf{La frase da avere pronta per l'esame:} \emph{``Lo stesso identico algoritmo Primal-Dual eccelle nel drift continuo di Req3 (dove raggiunge il 98\% circa di spesa budget) e resta strutturalmente più conservativo negli shock discreti di Req4 (91\% anche dopo retuning) -- non per un errore di implementazione, ma perché la garanzia di Hedge è di regret statico contro la miglior distribuzione fissa, non di tracking contro un target che cambia bruscamente.''}
\end{keybox}

%============================================================
\section{Domande d'esame plausibili}
%============================================================

\begin{qabox}{D: il professore si aspetta questo risultato "deludente" per Primal-Dual?}
Molto probabilmente sì. Il testo del requirement chiede esplicitamente di riusare l'algoritmo di Req3 come \emph{riferimento}, non di farlo vincere. Un risultato in cui Primal-Dual è sistematicamente più conservativo, con una spiegazione teorica precisa del \emph{perché}, è pedagogicamente più interessante di un pareggio a tre su tutto -- è la dimostrazione di aver capito il dominio di validità delle garanzie teoriche, non solo di saper implementare l'algoritmo.
\end{qabox}

\begin{qabox}{D: com'è collegato l'ambiente di Req4 a quello di Req3?}
Stessa classe \texttt{AdversarialMultiCampaignEnv}, stesso protocollo \texttt{round()}, stesso feedback completo su $m_t$. Cambia solo il metodo di generazione della sequenza: \texttt{\_build\_drift} (sinusoide continua per Req3) vs \texttt{\_build\_shocks} (blocchi piecewise per Req4). Tutti gli altri parametri del problema (N, B, T, valori, conflitti) sono identici e importati dallo stesso \texttt{req3\_config.py}.
\end{qabox}

\begin{qabox}{D: perché Sliding-Window è l'agente peggiore, e non semplicemente "diverso"?}
Con $W=500$ e blocchi da 2000 round, SW-UCB butta via dati validi molto prima che il blocco finisca, restando perennemente in una fase di esplorazione rumorosa invece di sfruttare stabilmente. \`E un compromesso deliberato (finestra corta = adattamento rapido ma varianza alta durante lo sfruttamento); con $W$ pari alla lunghezza del blocco l'agente adatterebbe più lentamente ma sfrutterebbe meglio -- un secondo trade-off esplorabile, non mostrato nei risultati principali.
\end{qabox}

\begin{qabox}{D: il fatto che i regimi possano ripetersi tra i blocchi è un bug?}
No, è una scelta di design dell'ambiente (\texttt{\_build\_shocks} assegna i regimi ai blocchi in modo i.i.d., non a rotazione). \`E onesto segnalarlo come limite noto se chiesto, ma non invalida né i risultati né il benchmark piecewise, che usa comunque le distribuzioni vere -- ripetute o no -- di ciascun blocco specifico.
\end{qabox}

%============================================================
\section{Numeri chiave da ricordare}
%============================================================

\begin{keybox}
\begin{itemize}
\item Ambiente: 5 blocchi $\times$ 2000 round, $T=10\,000$, $B=1600$, stessi parametri di Req2/Req3.
\item Regret piecewise finale: \textbf{Primal-Dual 1426} $\approx$ \textbf{CUSUM 1430} $\ll$ \textbf{Sliding-Window 1948}.
\item Budget speso: CUSUM e SW $\approx$ \textbf{99.9\%}, Primal-Dual \textbf{91.0\%} (era 74.6\% coi tassi di Req3, non ritarati).
\item Retuning Primal-Dual: \texttt{hedge\_eta} $0.015 \to 0.16$, \texttt{ogd\_eta} $0.017 \to 0.003$.
\item CUSUM: $h=8.0$, $\varepsilon=0.02$, media \textbf{145 reset} su 38 celle per trial.
\item Messaggio centrale: il gap di Primal-Dual è \textbf{strutturale} (regret statico vs tracking), non risolvibile con il solo tuning dei learning rate.
\end{itemize}
\end{keybox}

\end{document}
